A single overall leaderboard hides the most scientifically interesting fact about a model:\emph{where} it is good. This view documents the three machinery pieces that turn per-domain metrics into comparative claims, the Spearman domain-correlation matrix  and the seven literature-grounded capability profiles that convert those numbers into a named diagnosis.


\subsubsection{Domain correlation: redundancy vs complementarity}

The \code{domain\_correlation} plot is a model\(\times\)model Spearman \(\rho\) matrix built from each model's vector of per-domain success rates. Two models with \(\rho\)\(\approx\)1 succeed and fail on the same domains - they are \emph{redundant} probes of capability. Two with low or negative \(\rho\) stress different things - they are \emph{complementary}, and a benchmark wants both. It is a tool for curating the model set and the domain set, not for ranking.

\begin{lstlisting}
sr_pivot = df_metrics.pivot_table(index="Model", columns="Domain", values="Success_Rate")
corr = sr_pivot.T.corr(method="spearman")   # Model \(\times\) Model rank correlation
\end{lstlisting}

\begin{bmcavbox}{Correlation needs a domain spread}
Spearman over success-rate vectors is only meaningful with several domains and some variance in them. if In the bundled run the domain count is small and many models sit at SR\(\approx\)0,\(\rho\) is dominated by ties and should be read as illustrative of the method, not as a stable finding.
\end{bmcavbox}

\subsubsection{Breadth of competence: stacked PS and the SR heatmap}

\code{ps\_by\_domain\_stacked} stacks each model's per-domain Planning Score contributions into one bar. Two models can reach the same total bar height by opposite routes: a tall single segment (high PS concentrated in one domain) versus many even segments (competence spread across domains). The stack makes "broad vs concentrated" visible at a glance exactly the distinction the overall PS scalar destroys.

The \code{success\_rate\_heatmap} is the raw companion: model\(\times\)domain SR with no normalisation, so here unlike the rank heatmap cells \emph{are} comparable across the whole grid. Use it to spot empty columns (domains no model solved \(\rightarrow\) an infrastructure or difficulty problem, not a model problem) and empty rows (models that solved nothing).

\subsubsection{The seven capability profiles and How a model is assigned a profile}

The profiles are the taxonomy made formal: each is a conjunction of threshold conditions over
\(\{SR,\allowbreak FASR,\allowbreak IWSR,\allowbreak RG,\allowbreak IHR,\allowbreak ER,\allowbreak PAS,\allowbreak CoT,\allowbreak p_{\mathrm{valid},k,\mathrm{flat}}\}\),, and each maps to a specific claim in the planning literature. They are the bridge from "here are nine numbers" to "this model is doing X".

Assignment is deliberately conservative. Every profile is evaluated and the \textbf{first exact match}(all required conditions met, none unavailable) wins. If nothing matches exactly the model is labelled \code{Mixed/Unclassified} and reported against its \emph{closest} profile - the one with the highest matched-condition count, ties broken by fewest required conditions.

\begin{lstlisting}
evaluations = [evaluate_profile_conditions(metrics, p) for p in PROFILE_DEFINITIONS]
exact = next((e for e in evaluations if e["exact"]), None)
if exact is None:
    best = max(evaluations, key=lambda e: (e["match_count"], -e["required_count"]))
    assigned = "Mixed/Unclassified"      # honest about partial matches
else:
    assigned = exact["profile"]
\end{lstlisting}

\begin{table}[h]\small
\centering
\begin{tabularx}{\linewidth}{@{}l I P@{}}
\toprule
Profile & Conditions & Interpretation \(\cdot\) reference\\
\midrule
\textbf{Genuine Planner} & \(\mathrm{SR}>0.70 \cdot \mathrm{FASR}>0.50 \cdot \mathrm{RG}<0.10 \cdot \mathrm{IHR}>0.90 \cdot \mathrm{CoT}>0.70\) & Reads the schema, solves first try, reasoning tracks the plan. Strongest claim\\
\addlinespace
\textbf{Efficient Corrector} & \(\mathrm{SR}>0.50 \cdot \mathrm{FASR}<0.30 \cdot \lvert\mathrm{SR}-\mathrm{IWSR}\rvert\le 0.10 \cdot 0.10\le \mathrm{RG}\le 0.30\) & Rarely first-try, but repairs land early so IWSR stays near SR. Stechly 2023 (verify via P(Valid|k))\\
\addlinespace
\textbf{Stochastic Searcher} & \(\mathrm{SR}\ge 0.30 \cdot \mathrm{FASR}<0.20 \cdot \mathrm{RG}>0.30\) \(\cdot\) [opt] \(P(\mathrm{Valid}\mid k)\) flat & Success inflated by retries; attempts behave like resampling. Stechly 2023\\
\addlinespace
\textbf{Lucky Retriever} & \(\mathrm{SR}>0.60 \cdot \mathrm{FASR}>0.50 \cdot \mathrm{IHR}<0.30 \cdot \mathrm{CoT}<0.40\) & Solves from surface patterns; weak grounding, post-hoc reasoning. Kambhampati 2024 \(\cdot\) Valmeekam obfuscation\\
\addlinespace
\textbf{Understander} & \(\mathrm{SR}<0.20 \cdot \mathrm{FASR}<0.10 \cdot \mathrm{ER}>0.70 \cdot \mathrm{IHR}>0.90 \cdot \mathrm{PAS}>0.70\) & Knows vocabulary \&local preconditions, fails long-horizon order. Valmeekam Phase 2 \(\cdot\) Dziri 2023\\
\addlinespace
\textbf{Vocabulary-Only} & \(\mathrm{SR}<0.15 \cdot \mathrm{IHR}>0.90 \cdot 0.30\le \mathrm{ER}\le 0.70 \cdot \mathrm{PAS}<0.30\) & Legal action names, no coherent state-transition model. Kambhampati 2024 \(\cdot\) Vafa 2024\\
\addlinespace
\textbf{No Grounding} & \(\mathrm{SR}\le 0.05 \cdot \mathrm{IHR}<0.60 \cdot \mathrm{ER}<0.30 \cdot \mathrm{PAS}<0.30\) & Not reading the formal domain at all; everything ungrounded. Valmeekam 2023 \(\cdot\) McCoy 2024\\
\bottomrule
\end{tabularx}
\end{table}


\begin{bmkeybox}{NaN-aware conditions: "unavailable" \(\neq\) "missing"}
A condition whose metric is NaN (e.g. CoT when reasoning was never produced) is marked \emph{unavailable}, not \emph{failed}.
This is what stops a model being thrown out of "Genuine Planner" merely because CoT could not be computed - and it is why the report exposes the full matched/missing/unavailable breakdown per profile, so you can see exactly how close a model sat to each boundary rather than trusting a bare label.
\end{bmkeybox}

\begin{bmcavbox}{Profiles are thresholds, not ground truth}
The cut-points (\(\mathrm{SR}>0.70\), \(\mathrm{IHR}>0.90\), ...) encode the literature's qualitative claims as hard numbers; they are defensible but not sacred. A model just under a boundary is reported as Mixed against its closest profile precisely so the threshold choice stays visible and auditable. The label as to be treated as hypothesis that the condition breakdown lets you check, not a verdict.
\end{bmcavbox}

\subsubsection{Reading the metric set together: the diagnostic framework}

No single metric is a verdict. The scientific meaning lives in \emph{combinations} and the within-domain ranking heatmap is where they are read side by side. There are three reading directions, and then a small set of metric pairs that must never be read in isolation.

\paragraph{Three ways to read the ranking heatmap.}

\begin{table}[h]\small
\centering
\begin{tabularx}{\linewidth}{@{}l Y@{}}
\toprule
Direction & What it tells you\\
\midrule
A \textbf{row}(one model) & Scan all metrics: uniformly green = leads on every dimension here; uniformly red = worst on everything; a mixed row is a capability profile --- strengths and weaknesses at once.\\
\addlinespace
A \textbf{column}(one metric) & Which model dominates that metric. A column all one colour means the metric doesn't discriminate in this domain.\\
\addlinespace
\textbf{Cell blocks} & Green SR+FASR but red halluc+CoT = succeeding despite poor grounding (pattern-matching). Red SR but green exec+PAS = failing at the plan level yet showing process-level understanding.\\
\bottomrule
\end{tabularx}
\end{table}

\paragraph{The metric pairs to always read together.}

\textbf{Capability triangle --- SR + FASR + IWSR + RG}(is performance genuine or scaffolded by retrieval). Read SR and FASR first, then RG = SR - FASR, then IWSR position.

\begin{table}[h]\small
\centering
\begin{tabularx}{\linewidth}{@{}C C C C Y@{}}
\toprule
SR & FASR & RG & IWSR & Diagnosis\\
\midrule
High & High & \(\approx\)0--0.10 & \(\approx\)FASR & Genuine zero-shot planner --- retries irrelevant.\\
\addlinespace
High & Medium & 0.10--0.25 & \(\approx\)SR & Efficient corrector --- converges on attempts 2--3.\\
\addlinespace
High & Low & \(>0.30\) & \(\approx\)FASR & Stochastic search --- SR inflated; late retries are re-samples.\\
\addlinespace
Medium & Low & 0.20--0.40 & mid & Partial capability, retry-scaffolded.\\
\addlinespace
Low & Low & \(\approx\)0--0.05 & \(\approx\)FASR & No capability --- nothing for retries to inflate.\\
\bottomrule
\end{tabularx}
\end{table}

\textbf{Failure anatomy --- ER + PAS}(most useful when \(\mathrm{SR}<0.15\): how deep does it run, and why did it stop )

\begin{table}[h]\small
\centering
\begin{tabularx}{\linewidth}{@{}C C I P@{}}
\toprule
ER & PAS & What the model is doing & Valmeekam phase\\
\midrule
\(>0.70\) & \(>0.70\) & Runs deep, ordering errors near the goal --- understands domain, mis-sequences late. & Phase 2 (long-horizon)\\
\addlinespace
\(>0.70\) & \(<0.30\) & Runs deep but fabricates states --- knows vocabulary, invents transitions. & Phase 1/2 boundary\\
\addlinespace
\(<0.30\) & \(>0.70\) & Fails early, but on ordering of the few real steps. Rare. & Phase 2 + vocab issues\\
\addlinespace
\(<0.30\) & \(<0.30\) & Fails immediately, invents states from step 1. & Phase 1 (vocabulary)\\
\addlinespace
Medium & NaN & Mixed; PAS undefined ->check halluc gate and temporal distance. & Mixed\\
\bottomrule
\end{tabularx}
\end{table}

\textbf{Reasoning test SR + CoT} and \textbf{refinement pair IHR + PAS}(genuine vs retrieved success; and which kind of failure).

\begin{table}[H]\small
\centering
\begin{tabularx}{\linewidth}{@{}Y I@{}}
\toprule
Pair & Meaning\\
\midrule
\(\mathrm{SR}\) high \(\cdot\) \(\mathrm{CoT}>0.70\)& Reasoning guides grounded plans --- strongest evidence of planning (Lanham).\\
\addlinespace
\(\mathrm{SR}\) high \(\cdot\) \(\mathrm{CoT}<0.40\)& Fake reasoning --- correct plans, post-hoc CoT (Kambhampati).\\
\addlinespace
\(\mathrm{SR}\) low \(\cdot\) \(\mathrm{CoT}>0.70\)& Faithful but failing --- reasoning grounded, execution fails.\\
\addlinespace
\(\mathrm{IHR}>0.90\) \(\cdot\) \(\mathrm{PAS}>0.70\)& Knows vocabulary, ordering errors only --- most fixable.\\
\addlinespace
\(\mathrm{IHR}>0.90\) \(\cdot\) \(\mathrm{PAS}<0.30\)& Knows words, lacks world model.\\
\addlinespace
\(\mathrm{IHR}<0.60\) \(\cdot\) \(\mathrm{PAS}\) low& Near-random output --- no grounding, no causal model.\\
\bottomrule
\end{tabularx}
\end{table}
\vspace{-\parskip}
\begin{bmcavbox}{IHR is a gate, not just a metric}
If inverse-hallucination rate falls below \textasciitilde{} 0.6, ER, PAS and CoTaL are all unreliable: the model is failing at vocabulary, so the downstream causal metrics have nothing trustworthy to measure. Always clear the IHR gate before interpreting the process metrics.
\end{bmcavbox}

\paragraph{The diagnostic tree.}

The pairs above compose into a single ordered procedure read the metrics in this sequence, taking the branch that matches, and the model classifies itself.

{\small\noindent\begin{tabularx}{\linewidth}{@{}P I@{}}
\toprule
\(\mathrm{IHR}<0.60\) & STOP. Not reading the schema --- vocabulary failure. Do not interpret exec/PAS/CoT.\\
\addlinespace
IHR 0.60--0.90 & Proceed with caution; note partial vocabulary issues.\\
\addlinespace
\(\mathrm{IHR}>0.90\) & Proceed to Level 1.\\
\bottomrule
\end{tabularx}}

\smallskip{\small\noindent\begin{tabularx}{\linewidth}{@{}P I@{}}
\toprule
\(\mathrm{SR}<0.15\) & Almost no valid plans --- skip to Level 4; the question is which failure dominates.\\
\addlinespace
SR 0.15--0.50 & Partial capability --- Level 2.\\
\addlinespace
\(\mathrm{SR}>0.50\) & Can plan here --- Level 2 to test authenticity.\\
\bottomrule
\end{tabularx}}

\smallskip{\small\noindent\begin{tabularx}{\linewidth}{@{}P I@{}}
\toprule
\(\mathrm{RG}<0.10\) & Retry-independent (FASR\(\approx\)SR) --- Level 3 + Level 5.\\
\addlinespace
RG 0.10--0.30 & Moderate dependency --- report SR and FASR; Level 3.\\
\addlinespace
\(\mathrm{RG}>0.30\) & Strong dependency --- SR inflated, use FASR; confirm via flat P(Valid|k).\\
\bottomrule
\end{tabularx}}

\smallskip{\small\noindent\begin{tabularx}{\linewidth}{@{}P I@{}}
\toprule
IWSR \(\approx\) FASR & All success on attempt 1 --- confirm with declining P(Valid|k).\\
\addlinespace
IWSR near SR & Efficient corrector (attempts 2--3) --- verify with rising P(Valid|k).\\
\addlinespace
IWSR near FASR, high RG & Late-retry luck (attempts 4--5) --- stochastic; flat P(Valid|k).\\
\bottomrule
\end{tabularx}}

\smallskip{\small\noindent\begin{tabularx}{\linewidth}{@{}P I@{}}
\toprule
\(\mathrm{ER}>0.70\) \(\cdot\) \(\mathrm{PAS}>0.70\) & Long-horizon sequencing failure --- Phase 2 (Valmeekam).\\
\addlinespace
\(\mathrm{ER}>0.70\) \(\cdot\) \(\mathrm{PAS}<0.30\) & Runs deep but fabricates --- world-model incoherence (Vafa).\\
\addlinespace
\(\mathrm{ER}<0.30\) \(\cdot\) \(\mathrm{PAS}<0.30\) & Fails immediately --- Phase 1 fabrication.\\
\addlinespace
ER 0.30--0.70 & Mixed --- check temporal distance (low = near-miss, high = broken causal model).\\
\bottomrule
\end{tabularx}}

\smallskip{\small\noindent\begin{tabularx}{\linewidth}{@{}P I@{}}
\toprule
\(\mathrm{CoT}>0.70\) \(\cdot\) \(\mathrm{SR}\) high & Genuine deliberation --- strongest result.\\
\addlinespace
\(\mathrm{CoT}<0.40\) \(\cdot\) \(\mathrm{SR}\) high & Fake reasoning --- post-hoc CoT (Kambhampati); report as caveat.\\
\addlinespace
\(\mathrm{CoT}>0.70\) \(\cdot\) \(\mathrm{SR}\) low & Faithful but failing --- likely long-horizon (if ER also high).\\
\addlinespace
\(\mathrm{CoT}<0.40\) \(\cdot\) \(\mathrm{SR}\) low & Complete grounding failure --- return to Level 0 for a masked vocab issue.\\
\bottomrule
\end{tabularx}}
